\documentclass[technicalreport]{classes/ieicej}
% \usepackage[dvipdfmx]{graphicx,xcolor}
% %\usepackage[dvips]{graphicx}
\usepackage[fleqn]{amsmath}
%\usepackage{amsthm}
\usepackage{newtxtext}% 英数字フォントの設定を変更しないでください
\usepackage[varg]{newtxmath}% % 英数字フォントの設定を変更しないでください
%\usepackage{amssymb}
%\usepackage{bm}
\usepackage{subfiles}
\usepackage{url}
% ieicej.clsのciteコマンドとhyperrefは互換性がないため、hyperrefのcite関連機能を無効化
\usepackage[dvipdfmx,hidelinks,unicode,pdfusetitle,hyperfootnotes=false]{hyperref}
% hyperrefがbibciteを上書きしないようにする
\makeatletter
\let\hyper@bib@natlinkstart\@gobble
\let\hyper@bib@natlinkend\relax
\let\hyper@citelink@start\@gobble
\let\hyper@citelink@end\relax
\renewcommand{\hyper@natlinkstart}[1]{}%
\renewcommand{\hyper@natlinkend}{}%
\renewcommand{\hyper@natanchorstart}[1]{}%
\renewcommand{\hyper@natanchorend}{}%
% bibciteをhyperrefから保護
\let\orig@bibcite\bibcite
\renewcommand{\bibcite}[2]{\expandafter\gdef\csname b@#1\endcsname{#2}}
\makeatother

\setcounter{page}{1}

\field{}
\jtitle{実環境を見据えた音による行動認識モデルの精度向上}
\etitle{Improving Accuracy of Audio-Based Activity Recognition Models for Real-World Environments}
\authorlist{%
%  \authorentry{}{}{}\MembershipNumber{}
 \authorentry[s1f102201278@iniad.org]{竹本 志恩}{Shion Takemoto}{INIAD}\MembershipNumber{}
 \authorentry[]{川原 亮一}{Ryoichi Kawahara}{INIAD}\MembershipNumber{}
 %\authorentry[メールアドレス]{和文著者名}{英文著者名}{所属ラベル}\MembershipNumber{}
 %\authorentry{和文著者名}{英文著者名}{所属ラベル}[現在の所属ラベル]\MembershipNumber{}
}
% \affiliate[]{}{}
\affiliate[INIAD]{東洋大学 情報連携学部 情報連携学科}{Toyo Univ. INIAD?}
%\paffiliate[]{}
%\paffiliate[現在の所属ラベル]{和文所属}
\jalcdoi{???????????}% ← このままにしておいてください

\begin{document}
\begin{abstract}
プライバシーに配慮した高齢者見守りに向け，音響イベント検出の実環境適用を目指す．
本研究では，モデルパラメータを増大させずに精度を向上させるため，信頼度を用いた半教師あり学習を改良し，ドメイン汎化手法のアーキテクチャを検討した後，後処理手法含む三手法を統合した．
実験の結果，モデル規模を維持したまま先行研究を上回る精度を達成し，提案手法の有効性を確認した．
\end{abstract}

\begin{keyword}
%和文キーワード 4〜5語
\end{keyword}
\begin{eabstract}
%英文アブストラクト 100 words
\end{eabstract}
\begin{ekeyword}
%英文キーワード
\end{ekeyword}
\maketitle

% These commands are not defined in ieicej.cls - use \jtitle, \etitle, and \authorlist instead
%\title{実環境を見据えた\\ 音による行動認識モデルの精度向上}
%\author{竹本 志恩 }
%\studentid{1F10220127}
%\supervisor{川原 亮一 教授}
% \date{2025年 11月 30日}

\section{まえがき}
% \subfile{sections/00_abstract.tex}
\subfile{sections/01_introduction.tex}
% \subfile{sections/02_main_issue.tex}
% \subfile{sections/03_eval_exp.tex}
% \subfile{sections/04_eval_res.tex}
% \subfile{sections/05_discussion.tex}


\ack %% 謝辞

\bibliographystyle{sieicej}
\bibliography{references}


\appendix
% \section{}
% \subfile{sections/appendix.tex}

%% 著者紹介・顔写真の掲載はC分冊の場合は任意です．
\begin{biography}
\profile*{正会員}{竹本 志恩}{} %学生会員か不明
%\profile{会員種別}{名前}{紹介文}% 顔写真あり
%\profile*{会員種別}{名前}{紹介文}% 顔写真なし
\end{biography}

\end{document}